\section{Compensation in a fixed integral homology class}\label{p14-sec:compact}
Let \((X,\omega)\) be a closed symplectic six-manifold with
\(c_1(TX)=0\). At an endpoint \(J\) of a smooth path of
\(\omega\)-tame structures, mark finitely many disjoint embedded
curves \(C_i\) in nonzero integral classes \(A_i\). Suppose each admits
arbitrarily small tubes in which the only cycles below a prescribed energy
bound are its multiples. Suppose also that each marked cycle series
has a geometric factorization into elementary clusters. These
hypotheses hold for the super-rigid homogeneous germs used in
Theorem~\ref{p14-thm:wall}: projection expresses each irreducible compact
curve in the bundle as a finite map with a height section, and a
vertical fiber contains no compact curve. The cluster construction
below then applies to the insertion consisting of the multiple zero
sections at each energy bound.

Let \(H\) be the virtual-dimension-zero, Chern-number-zero part of
Pardon's rational algebra of complex curves \cite{pardon}.
Write \(S\) for its almost-complex cycle-space counterpart and
\(\Phi:H\to S\) for the comparison of \cite[\S2]{WueProduct}. In \(S\)
a compact open and closed part of the relative cycle space defines an
insertion. Excision on disjoint supports gives multiplication, and
the two restrictions of such an insertion over a closed interval
agree. Denote by \(F_h\) the normalized balanced factor of a genus-\(h\)
elementary curve: it is the equivariant series on the sum of two
degree-\(h-1\) line bundles with product the canonical bundle and
circle weights \(+1,-1\). Its linear term is the multiplicity-one
class \(x_{h,1}\) and its image under \(\Phi\) is the elementary
cluster series \cite[\S\S3--4]{WueProduct}. We use the full integral
label \(\beta\in H_2(X;\ZZ)\) in the monomial \(Q^\beta\).
Products at a fixed energy cutoff are taken modulo terms above that
cutoff. For each marked curve send the local variable \(q^d\) to
\(Q^{dA_i}\), and let \(P_J\) be the product of its marked factors
over all the marked curves.

\begin{proposition}\label{p14-prop:exterior}
At each such endpoint there is a product \(R_J\) of elementary factors
with integral exponents, obtained from the connected supports outside
the marked sectors after a generic perturbation. It is independent
of the perturbation, the cluster decomposition, and sufficiently
small isolating tubes. For endpoints of a smooth tame path,
\begin{equation}\label{p14-eq:conservation}
                    P_{J_-}R_{J_-}=P_{J_+}R_{J_+}.
\end{equation}
The equality holds in the class-labelled completion with coefficients
in \(H\). It holds also when the marked curves differ at the two ends.
\end{proposition}
\begin{proof}
First take an integral taming form and fix an energy cutoff \(L\).
Choose smaller disjoint tubes \(V_i\) with closure in the respective
isolation neighborhoods. At the endpoint no nonempty cycle contained
in \(\overline V_i\) meets its boundary. This persists for nearby
structures: a contrary sequence would have a bounded-mass subsequence
converging as cycles with Hausdorff convergence of supports
\cite[Theorem~2.11]{diw}. A nonempty limit has positive energy, and
integrality prevents a sequence of positive cycles from limiting to
the empty cycle. Thus containment in each tube is open and closed
in the bounded-energy cycle space over a short path to a generic
endpoint perturbation.

At the perturbed endpoint distinct irreducible supports are disjoint.
The bounded-energy space of connected supports, counted with possible
multiplicities in the compactification, splits into compact open and
closed parts according to containment in the marked tubes and their
complement. To see closedness, a limit of connected supports is
connected, and disjointness of distinct irreducible curves forces it to
have a single irreducible support. The respective unions of supports
are disjoint compact subsets of \(X\). Restriction to disjoint open
neighborhoods splits each effective current into currents labelled by
the parts. This factors the entire insertion in \(S\), with mixed
configurations included by multiplying the part series. In particular
\(R_J\) is constructed from configurations all of whose connected
supports are complementary.

Subdivide each part by a finite cluster partition. Cluster isotopy,
excision and the fold relation decompose each cluster into a signed
elementary factor and corrections in classes \(d\gamma\), \(d\geq2\),
if the original class is \(\gamma\)
\cite[Propositions~3.18, 3.21 and Theorem~3.25]{diw}
\cite[\S4]{WueProduct}. The latter reference gives the identities of
cycle insertions needed here. Each step decreases
\(\lfloor L/\omega(\gamma)\rfloor\). This yields finite products for
the parts with integral exponents. The image of their product under
\(\Phi\) is precisely the insertion of the entire target at this
cutoff. For each integral class the entire class locus is compact
over the original path by cycle compactness; the class is locally
constant. The interval relation identifies its two endpoint
insertions. The argument requires isolation of the marked sectors
only at the endpoints.

To prove uniqueness, cancel factors in order of increasing positive
energy. At a first differing label \(\beta\), the difference of
the comparison images is a finite sum
\(\sum_h d_h\Phi(x_{h,1})\). All nonlinear terms at this label are
composed of smaller-energy contributions. Other labels at the same
energy have distinct monomials, also when they differ by torsion.
Independence of the multiplicity-one classes
\cite[Theorem~5.1]{WueProduct} gives \(d_h=0\) for every \(h\). This
proves equality of the two total products, uniqueness at an endpoint
and consistency on increasing the cutoff. Applying the same argument
in each isolating region fixes the marked factors. Equation
\eqref{p14-eq:conservation} and independence of the exterior follow.

For an arbitrary taming form, approximate its cohomology class on
the compact path by a rational class represented by a nearby closed
taming form and rescale the latter to integral periods. The energies
are uniformly comparable on the complex lines along this path. At a
fixed endpoint let \(M_J\) be the monoid of classes of effective cycles.
Any nonzero exponent of the full product has label in \(M_J\):
induct on positive integral energy. At an ineffective class the entire
insertion is zero and smaller effective labels have sums in \(M_J\).
Multiplicity-one detection forces all the new exponents there to
vanish. This also applies to the marked products in their tubes.
Cycle compactness and local constancy of class imply finiteness of
effective labels at bounded energy. Two choices of taming form
therefore give products in a common locally finite completion at this
endpoint, since their energies are comparable on \(M_J\).
Their images coincide at each label and the same induction identifies
their exponents. Subtracting the marked multiplicities identifies the
complementary ones as well.
\end{proof}

For example, in class $7A$ over the genus-two core with two branch
values the change in the marked Euler integral of source genus
thirteen is
\[
 a_{\{1,2\}}-\tfrac{11}{12}(a_{\{1\}}+a_{\{2\}}).
\]
Indeed $x^{11}=v^{22}-(11/12)v^{24}+O(v^{26})$, whereas
$x^{12}=v^{24}+O(v^{26})$. In the elementary product the two genera
give separate integral changes. The exterior multiplicities in
genera twelve and thirteen have the opposite changes, respectively,
by Proposition~\ref{p14-prop:exterior}.

\begin{proof}[Proof of Theorem~\ref{p14-thm:torus}]
On \(X=\RR^6/\ZZ^6\) with its standard symplectic form there exists a
constant tame linear complex structure admitting no nonempty compact
effective curve. To see this, identify the integral classes with
integral bivectors. For every nonzero integral bivector \(B\) the
equation \((\Lambda^2J)B=B\) defines a proper analytic subset of the
space of positive linear complex structures. Indeed this space is the
connected homogeneous space \(\mathrm{GL}^+(6,\RR)/\mathrm{GL}(3,\CC)\).
If the equation held on an open subset, the standard \(J\) would fix
the span of the entire \(\mathrm{GL}^+\)-orbit of \(B\) in the exterior
square. A diagonal matrix with distinct pairwise products of its
diagonal entries extracts the nonzero coordinate bivectors by
polynomial interpolation; oriented signed permutations then give every
basis bivector in the span. The standard \(J\) does not fix a wedge
of real coordinate vectors from two different complex lines.
Baire's theorem applied to the open tame subset gives a constant
\(J\) avoiding these equations for all \(B\ne0\). The structure descends
under translations. Every effective current would give an invariant
integral bivector by integration of its complex tangent planes, with
positive symplectic area. Hence the total series is one there.
The space of tame structures is path connected, so
Proposition~\ref{p14-prop:exterior} gives \(R_J=P_J^{-1}\) at every marked
endpoint on this torus.

In a product of two symplectic two-tori smooth the positive transverse
intersections of one horizontal torus with \(g-1\) vertical tori,
using \(zw=\delta\) in disjoint coordinate balls. The intersection
graph is a tree. The result is a genus-\(g\) symplectic surface of
class \([T_1]+(g-1)[T_2]\). Take its product with a point of the last
two-torus. The resulting class \(A\) is primitive by pairing with
the area class of the first factor. Its symplectic normal bundle has
Chern number \(2g-2\); unitary bundles of rank two over a surface are
classified by this number. A tubular identification thus realizes the
given homogeneous structure with its prescribed normal derivative.
It is tame near the zero section. Cutting off in a chart of tame
structures extends it over \(X\) while preserving the germ. At every
positive even \(b\) a connected cubic as required exists on the same
base: send all puncture loops to \((12)\), one handle generator to
\((123)\), and the rest to the identity, and fill the finite cover
over the punctures. Applying Theorem~\ref{p14-thm:wall} changes the
exterior multiplicities by \(-c_j\) at \(h_j\), and by zero in every
other genus at label \(7A\). The binomial bound of
\S\ref{p14-sec:chambers} proves sharpness.
\end{proof}

In particular at least one of the two endpoints attaining this
maximum has \(\sum_h|n^{\rm ext}_{7A,h}|\geq2^{b-1}\), by the
triangle inequality and integrality. The area and the class are fixed
as \(b\) increases. These endpoints may approach a smooth reference
structure whose marked section is itself super-rigid. Use, for example,
the homogeneous split spin model of \cite[Lemma~14.1]{WueUnramified} on the
same bundle. All the homogeneous structures agree with it along
\(C\). On a tube of radius \(R\), cutting off their difference in a
smooth chart gives sup norm at most \(K_b R\) and first derivative
bounded by \(K_b\); the constant may depend on the crossing. For each
\(b\) take \(R_b\) so small that \(K_b R_b^{1/b}\leq 1/b\).
Interpolation yields convergence in every \(C^{0,\alpha}\),
\(0<\alpha<1\). Real dilations of the full-join graphs in
\S\ref{p14-sec:smooth} likewise place their embedded representatives
in arbitrarily small tubes, using the operators of that construction.
Their genus is \(h_b=1+7(g-1)+5b/2\). It is unbounded along this
sequence of smooth tame structures at fixed symplectic area.

\smallskip
The compensation determines the total signed change outside the
isolating tube at the two endpoints. Allocating that change to the
individual intermediate curves of an arbitrary compact deformation
remains open. For simultaneous ramified crossings there is a further
local question: several independent heights can match on a common
component, and the resulting strata need not have a common
dimension or stabilizer. The single adjoint equation of
\S\ref{p14-sec:joins} treats the prescribed cubic crossing.
